% Import der Config-Variablen
%% -- Allgemeine Konfiguration -- %%

%Textgröße
\newcommand{\mainsize}{12pt}


%% -- Deckblattkonfiguration -- %%

%Name der Universität/Hochschule
\newcommand{\universityname}{Universitätsname}

%Name der Fakultät
\newcommand{\facultyname}{Fakultätsname}

%Dateiname des Logos
\newcommand{\logofilename}{logo}

%Arbeitstyp
\newcommand{\thesistype}{Bachelorarbeit}	

%Titel der Arbeit
\newcommand{\thesistitle}{Titel}

%Zielabschluss
\newcommand{\studygoal}{Bachelor of Science}

%Autor der Arbeit
\newcommand{\thesisauthor}{Autor}

%Studiengang
\newcommand{\studytype}{Studiengang}

%Matrikelnummer
\newcommand{\studentnumber}{Matrikelnummer}

%Semesternummer
\newcommand{\semesternumber}{Fachsemester}

%Mail
\newcommand{\mail}{.edu-Mail}

%Abgabetermin
\newcommand{\submissiondate}{Abgabetermin}

%Semester
\newcommand{\semester}{WS 2020}

%Betreuer (mit akademischen Titeln)
\newcommand{\supervisor}{Betreuer}


%% -- Inhaltsverzeichnis - Konfiguration -- %%

%Einzug Subsection
\newcommand{\idxsubsecindent}{0.25cm}

%Einzug Subsubsection
\newcommand{\idxsubsubsecindent}{0.6cm}

% Import der Befehle
% Befehl um Textbox für den Druck zu generieren
\newcommand{\calibrationbox}[2]{% parameters: #1=width, #2=height
	\setlength{\unitlength}{1.0mm}%
	\begin{picture}(#1,#2)%
	\linethickness{0.05mm}%
	\put(0,0){\dashbox{0.2}(#1,#2)%
	{\parbox{#1mm}{%
	\centering\footnotesize 
	width $ = #1 \textrm{mm}$\\
	height $ = #2 \textrm{mm}$
	}}}\end{picture}
}

% Eigener Befehl für autoref + Seitenangabe
\newcommand{\autorefpage}[1]{\autoref{#1} (S. \pageref{#1})}

% Befehl für Quellenangaben direkt unter dem Bild (rechts)
\newcommand*{\sourceright}[1]{\par\raggedleft\footnotesize \textit{Quelle}:~#1}

% Befehl für Quellenangaben direkt unter dem Bild (zentriert)
\newcommand*{\sourcecenter}[1]{\par\centering\footnotesize \textit{Quelle}:~#1}

% Metadatenangaben
\author{\thesisauthor}
\title{\thesistitle}

\documentclass[\maintextsize, \papersize, fleqn, xcolor=dvipsnames]{scrartcl}

% PDF-Inhalt such- und kopierbar machen
\usepackage[resetfonts]{cmap}

% Anpassungen für deutsche Sprache
\usepackage[utf8]{inputenc}
\usepackage[T1]{fontenc}
\usepackage[ngerman]{babel}

% Mehr Farben zulassen
\usepackage[table,dvipsnames]{xcolor}

% Grafiken
\usepackage{graphicx}
\graphicspath{{images/}}
\usepackage{float}

% Textschriftart ändern
\usepackage{lmodern}

% Verbessertes Fontrendering
\usepackage[final]{microtype}

% Seitenränder festlegen
\usepackage[\papersize,left=\leftpapermargin,right=\rightpapermargin,top=\toppapermargin,bottom=\bottompapermargin]{geometry}

% 1,5 Zeilenabstand konfigurieren
\usepackage[onehalfspacing]{setspace}

% Zeileneinschub mit Abstand ersetzen
\usepackage{parskip}

% Fußnoten konfigurieren
\usepackage[perpage, flushmargin, hang]{footmisc}

% Inhaltsverzeichnis Aussehen konfigurieren und ins Inhaltsverzeichnis mit aufnehmen
\usepackage{tocbibind}
\usepackage{tocloft}
\renewcommand{\cftsecleader}{\cftdotfill{\cftdotsep}}
\setlength{\cftsubsecindent}{\idxsubsecindent}
\setlength{\cftsubsubsecindent}{\idxsubsubsecindent}
\setlength{\cftbeforesecskip}{3pt}

% Schönere Seitenangaben
\usepackage{fancyhdr}
\fancyhf{}
\fancyhead[EL,OR]{\thepage}
\fancyhead[ER]{\chaptername~\thechapter}
\fancyhead[LO]{\nouppercase{\leftmark}}

% Namensnennungen kursiv schreiben
\usepackage{xpatch}
\xpatchcmd{\citeauthor}{\begingroup}{\begingroup\em}{}{}

% Mathematik-Pakete
\usepackage{amsmath,amsthm,amssymb}
\setlength{\mathindent}{1cm}

% Automatisches Setzen von Anführungszeichen (\enquote, \foreignquote, \blockquote)
\usepackage[autostyle=true,german=quotes]{csquotes}

% Einbinden von PDFs
\usepackage{pdfpages}

% Beschriftung von Bildern/Tabellen/Listings usw.
\usepackage[
	format=plain,
	margin=0.5cm,
	labelformat=simple,
	labelfont={bf},
	font=footnotesize
]{caption}

% Erlaubt die Erstellung von zusammengesetzten Grafiken
\usepackage{subcaption}

% Leere Anfangs- und Endzeilen vermeiden
\usepackage[defaultlines=4,all]{nowidow}

% Bessere Aufzählungen
\usepackage{paralist}

% Bessere Tabellen
\usepackage{booktabs}

% Euro-Zeichen
\usepackage[official]{eurosym}

% Abkürzungen
\usepackage[footnote, printonlyused, smaller]{acronym}

%%% === LISTINGS === %%%
\usepackage{listings}
\usepackage{beramono}

% CSS als Sprache definieren
\lstdefinelanguage{CSS}{
  keywords={color,background-image:,margin,padding,font,weight,display,position,top,left,right,bottom,list,style,border,size,white,space,min,width, transition:, transform:, transition-property, transition-duration, transition-timing-function},	
  sensitive=true,
  morecomment=[l]{//},
  morecomment=[s]{/*}{*/},
  morestring=[b]',
  morestring=[b]",
  alsoletter={:},
  alsodigit={-}
}

% JavaScript als Sprache definieren
\lstdefinelanguage{JavaScript}{
  keywords={break, case, catch, continue, debugger, default, delete, do, else, finally, for, function, if, in, instanceof, new, return, null, true, false, switch, this, throw, try, typeof, var, void, while, with},
  ndkeywords={class, export, boolean, throw, implements, import, this},
  sensitive=true,
  morecomment=[l]{//},
  morecomment=[s]{/*}{*/},
  morestring=[b]',
  morestring=[b]"
}

% HTML5 als Sprache definieren
\lstdefinelanguage{HTML5}{
  language=html,
  sensitive=true,	
  alsoletter={<>=-},	
  morecomment=[s]{<!-}{-->},
  tag=[s],
  otherkeywords={
  % General
  >,
  % Standard tags
	<!DOCTYPE,
  </html, <html, <head, <title, </title, <style, </style, <link, </head, <meta, />,
	% body
	</body, <body,
	% Divs
	</div, <div, </div>,
	% headings
	</h1, <h1, </h1>,
	</h2, <h2, </h2>, 
	</h3, <h3, </h3>, 
	</h4, <h4, </h4>, 
	</h5, <h5, </h5>,
	</h6, <h6, </h6>,  
	% Paragraphs
	</p, <p, </p>,
	% scripts
	</script, <script,
  % More tags...
  <canvas, /canvas>, <svg, <rect, <animateTransform, </rect>, </svg>, <video, <source, <iframe, </iframe>, </video>, <image, </image>, <header, </header, <article, </article
  },
  ndkeywords={
  % General
  =,
  % HTML attributes
  charset=, src=, id=, width=, height=, style=, type=, rel=, href=,
  % SVG attributes
  fill=, attributeName=, begin=, dur=, from=, to=, poster=, controls=, x=, y=, repeatCount=, xlink:href=,
  % properties
  margin:, padding:, background-image:, border:, top:, left:, position:, width:, height:, margin-top:, margin-bottom:, font-size:, line-height:,
	% CSS3 properties
  transform:, -moz-transform:, -webkit-transform:,
  animation:, -webkit-animation:,
  transition:,  transition-duration:, transition-property:, transition-timing-function:,
  }
}

% Kombination aus HTML,CSS,JS als Style definieren
\lstdefinestyle{htmlcssjs} {%
  % Code
  language=HTML5,
  alsolanguage=JavaScript,
  alsodigit={.:;},	
}

\lstset{
	breaklines=true,
	breakatwhitespace=false,
	captionpos=b,
	frame=shadowbox,
	columns=flexible,
	keepspaces=false,
	numbers=left,
	showspaces=false,
	showstringspaces=false,
	showtabs=false,
	identifierstyle=\color{black},
	ndkeywordstyle=\color{MidnightBlue}\bfseries,
	basicstyle=\footnotesize\ttfamily,
	numberstyle=\color{Gray},
	keywordstyle=\bfseries\color{RoyalBlue},
	commentstyle=\color{ForestGreen},
	stringstyle=\color{Orange},
	xleftmargin={0.75cm},
	tabsize=2,
	extendedchars=true,
	inputencoding=utf8,
	resetmargins=true,
	rulesepcolor=\color{lightgray},
	literate=
  {á}{{\'a}}1 {é}{{\'e}}1 {í}{{\'i}}1 {ó}{{\'o}}1 {ú}{{\'u}}1
  {Á}{{\'A}}1 {É}{{\'E}}1 {Í}{{\'I}}1 {Ó}{{\'O}}1 {Ú}{{\'U}}1
  {à}{{\`a}}1 {è}{{\`e}}1 {ì}{{\`i}}1 {ò}{{\`o}}1 {ù}{{\`u}}1
  {À}{{\`A}}1 {È}{{\'E}}1 {Ì}{{\`I}}1 {Ò}{{\`O}}1 {Ù}{{\`U}}1
  {ä}{{\"a}}1 {ë}{{\"e}}1 {ï}{{\"i}}1 {ö}{{\"o}}1 {ü}{{\"u}}1
  {Ä}{{\"A}}1 {Ë}{{\"E}}1 {Ï}{{\"I}}1 {Ö}{{\"O}}1 {Ü}{{\"U}}1
  {â}{{\^a}}1 {ê}{{\^e}}1 {î}{{\^i}}1 {ô}{{\^o}}1 {û}{{\^u}}1
  {Â}{{\^A}}1 {Ê}{{\^E}}1 {Î}{{\^I}}1 {Ô}{{\^O}}1 {Û}{{\^U}}1
  {Ã}{{\~A}}1 {ã}{{\~a}}1 {Õ}{{\~O}}1 {õ}{{\~o}}1
  {œ}{{\oe}}1 {Œ}{{\OE}}1 {æ}{{\ae}}1 {Æ}{{\AE}}1 {ß}{{\ss}}1
  {ű}{{\H{u}}}1 {Ű}{{\H{U}}}1 {ő}{{\H{o}}}1 {Ő}{{\H{O}}}1
  {ç}{{\c c}}1 {Ç}{{\c C}}1 {ø}{{\o}}1 {å}{{\r a}}1 {Å}{{\r A}}1
  {€}{{\euro}}1 {£}{{\pounds}}1 {«}{{\guillemotleft}}1
  {»}{{\guillemotright}}1 {ñ}{{\~n}}1 {Ñ}{{\~N}}1 {¿}{{?`}}1
}

%%% === ENDE LISTINGS === %%%

% Verlinkungen und PDF Daten initialisieren
\usepackage{hyperref}
\hypersetup{
	pdftitle = {\thesistitle},
	pdfauthor = {\thesisauthor},
	pdfcreator = {pdflatex},
	pdfkeywords={\thesistype, \thesistitle},
	pdfsubject={\thesistype},
	pdfproducer = {LaTeX with hyperref},
	colorlinks = false,
	linktoc = all,
	linkbordercolor = {0 0 0.5},
}
\usepackage[all]{hypcap}

%%% === ENDE DER KONFIGURATION === %%%



%%% === BEGINN DES DOKUMENTS === %%%

\begin{document}

%Namen der autoref generierten Verweise ändern
\renewcommand{\sectionautorefname}{Kapitel}
\renewcommand{\subsectionautorefname}{Abschnitt}
\renewcommand{\subsubsectionautorefname}{Unterabschnitt}

%Seitennummerierung abschalten
\pagestyle{empty}
	
%Deckblatt einbinden
%Seitenränder abschalten
\newgeometry{left=\titlemargin,right=\titlemargin,top=\titlemargin,bottom=\titlemargin}

\begin{titlepage}

   \singlespacing

   %Inhalt zentrieren
   \begin{center} \large 
    
    \universityname
    \vspace*{0.5cm}
    
    \facultyname
    \vspace*{1.5cm}
    
	\includegraphics[scale=0.45]{titlepage/\logofilename}
    \vspace*{2cm} 

	\thesistype \ zum Thema
	\vspace*{0.5cm}
	
    {\huge\textbf{\thesistitle}}
    \vspace*{1.5cm}
    
    
  	Zur Erlangung des Grades \studygoal
  	\vspace*{1.5cm}

    \textbf{\thesisauthor}
    \vspace*{0.5cm}
    
    \begin{tabular}{rl}
    	Studiengang: & \studytype \\
    	Matrikelnummer: & \studentnumber \\
    	Fachsemester: & \semesternumber \\
    	E-Mail: & \mail \\
    \end{tabular}
    
    \vspace*{2.0cm}

    Abgabe am: \submissiondate \\
    (\semester)
    \vspace*{1.5cm}

    Betreuung durch: \supervisor
  \end{center}
  
\end{titlepage}

%Seitenränder restaurieren
\restoregeometry
\newpage
	
%Leere Seite einfügen (Hier könnte z.B die handschriftliche Widmung stehen)
\mbox{}
\newpage
	
%Seitenangaben aktivieren
\pagestyle{fancy}
	
%Vorverzeichnisse einbinden
\begingroup

% Header auch für Seiten ohne Kopfzeile einfügen
\fancypagestyle{plain}

% Zeilenabstand ändern
\singlespacing

% Inhaltsverzeichnis einfügen
\tableofcontents
\newpage

% Abbildungsverzeichnis einfügen
\listoffigures
\newpage

% Tabellenverzeichnis einfügen
\listoftables
\newpage

% Quellcodeverzeichnis
\renewcommand{\lstlistoflistings}{\begingroup
\tocfile{\lstlistingname}{lol}
\endgroup}
\renewcommand{\lstlistingname}{Quellcodeverzeichnis}
\lstlistoflistings
\newpage

% Abkürzungsverzeichnis einfügen
\section*{Abkürzungsverzeichnis}
\markboth{Abkürzungsverzeichnis}{Abkürzungsverzeichnis}
\phantomsection
\addcontentsline{toc}{section}{\protect Abkürzungsverzeichnis}

\begin{acronym}
	\itemsep-20pt
	% Hier Abkürzungen eintragen mit \acro{Kürzel}[Kurzform]{Langform} bzw. \acroplural{Kürzel}[Kurzform]{Langform}
	% Wird die Pluralform benötigt, so muss unter dem selben Kürzel auch die Singularform hinterlegt sein!
	% Im Text kann dann \ac{Kürzel} bzw. \acp{Kürzel} verwendet werden.
	% =============================================================================================================
	
	\acro{html}[HTML]{Hypertext Markup Language}
	\acro{dr}[Dr.]{Doktor}
	\acroplural{dr}[Dres.]{Doktoren}
	
\end{acronym}
\newpage

% Symbolverzeichnis erstellen
\section*{Symbolverzeichnis}
\markboth{Symbolverzeichnis}{Symbolverzeichnis}
\phantomsection
\addcontentsline{toc}{section}{\protect Symbolverzeichnis}

\begin{acronym}
	\itemsep-20pt
	% Hier Abkürzungen eintragen mit \acro{Kürzel}[Kurzform]{Langform} bzw. \acroplural{Kürzel}[Kurzform]{Langform}
	% Wird die Pluralform benötigt, so muss unter dem selben Kürzel auch die Singularform hinterlegt sein!
	% Im Text kann dann \ac{Kürzel} bzw. \acp{Kürzel} verwendet werden.
	% =============================================================================================================
	\acro{a}[\ensuremath{A}]{{\acrounit{\meter^2}Oberfläche}}
\end{acronym}
\newpage

% Formelverzeichnis erstellen
\newlistof{equations}{equ}{Formelverzeichnis}
\renewcommand{\listofequations}{\begingroup
\tocfile{Formelverzeichnis}{equ}
\endgroup}
\listofequations
\newpage

\endgroup

%Vortexte einbinden
% === Einfügen eines Mottos, Sprichtworts, Zitat einer Persönlichkeit
%\pagestyle{empty}
\newgeometry{left=\leftpapermargin,right=\rightpapermargin,top=\titlemargin,bottom=\titlemargin}
\vspace*{\fill}
\begin{center}
	\huge{
	\glqq
	\textit{Das ist ein ganz wunderbares Zitat und beschreibt diese Arbeit sehr gut.}\grqq\\
	\textbf{--- Albert Einstein ---}
	}
\end{center}

\vfill


% === Einfügen eines Geleitworts
%\markboth{Geleitwort}{Geleitwort}
%\phantomsection
%\addcontentsline{toc}{section}{\protect Geleitwort}
%\section*{Geleitwort}\label{sec:geleitwort}

Das Geleitwort wird vom Herausgeber, dem Doktorvater oder einer Fachautorität verfasst.
Es umfasst inhaltlich entweder die Zuordnung der Arbeit in einen größeren Reihen- oder Fachzusammenhang oder aber eine kurze wertende Stellungnahme zu Gehalt und Bedeutung der Arbeit.

\newpage


% === Einfügen eines Vorworts
%\markboth{\preamblename}{\preamblename}
%\phantomsection
%\addcontentsline{toc}{section}{\protect \preamblename}
%% Dieses Dokument dient als Vorwort der Arbeit (wird aber meistens nur bei zu veröffentlichenden Arbeiten benötigt)
\section*{\preamblename}\label{sec:\preamblename}

% Beginn

Hier kann zum Beispiel auf den Anlass und die Bedeutung der Schrift und deren Stellung in der wissenschaftlichen Diskussion eingegangen werden.
Außerdem kann auf die sachliche und/oder finanzielle Unterstützung durch Dritte eingegangen werden.
Des Weiteren sollte auf Umfang und Quelle der erteilten Druckerlaubnisse hingewiesen und Dank ausgesprochen werden.

Art und Umfang der persönlichen Mitarbeit sowie Unterstützung durch Dritte ist zu erwähnen.

% Ende

\vspace{0.5cm}
\begingroup
\raggedright{\preamblelocation, \preambledate}
\hfill
\raggedleft{\textit{\thesisauthor}}
\endgroup

\newpage


% === Einfügen der Kurzzusammenfassung (alternativ in preindex)
%\markboth{Kurzzusammenfassung}{Kurzzusammenfassung}
%\phantomsection
%\addcontentsline{toc}{section}{\protect Kurzzusammenfassung}
%\section*{Kurzzusammenfassung}\label{sec:kurzzusammenfassung}

Hier würde normalerweise der eigentliche Text stehen. Folgende Kriterien sind einzuhalten:

 \begin{enumerate}
 	\item maximal eine Seite lang 150 - 250 Wörter und enthält eine kompakte Zusammenfassung zentraler Thesen und Ergebnisse.
 	\item Die eigene Vorgehensweise wird vorgestellt
 	\item Wenig bezug auf Literatur 
 	\item Kein Verweis auf offene Fragen und keine subjektive Bewertung o. Rechtfertigung
 	\item Zeitform: Präsens; vergangene Ereignisse im Perfekt
 \end{enumerate}

\newpage


% === Einfügen des Abstract (alternativ in preindex)
%\markboth{Abstract}{Abstract}
%\phantomsection
%\addcontentsline{toc}{section}{\protect Abstract}
%\section*{Abstract}\label{sec:abstract}

\selectlanguage{english}

This page has the same purpose as 'Kurzzusammenfassung'. The only difference is that it is written in english instead of german. The same criterias should be applied to this section.

\selectlanguage{ngerman}
\newpage

%Testtext
\section{Beispiele}\label{sec:beispiel}
	
Lorenawdawdwadz awod awd aiwd awdu awoud awduioawud oiawudoawuoduawoduawoidiuawduawodu awd awd awudwad uawdawiduawdawd adwaddwda\footnote{Testfußnote}

\subsection{Formeln}\label{subsec:formulas}

\begin{equation}\label{equ:relativitätstheorie}
	x^2 + mx + t = 0
\end{equation}
\equationentry{Geradengleichung}

Siehe \autoref{equ:relativitätstheorie}
	
\subsection{Bilder/Abbildungen}\label{subsec:abbildungen}
	
Einzelnes Bild
	
\begin{figure}[H]
	\centering
	\begin{minipage}{0.6\textwidth}
		\centering
		\fbox{\includegraphics[width=\textwidth]{logo}}
		\sourceright{Hier cite einfügen}
	\end{minipage}
	\caption{Logo der Hochschule München}
  	\label{fig:logo}
\end{figure}
	
Subfigures siehe \autorefpage{fig:coffee}
	
\begin{figure}[H]
  	\centering
  	\begin{subfigure}[b]{0.4\textwidth}
    	\includegraphics[width=\textwidth]{logo}
    	\caption{A logo.}
  	\end{subfigure}
  	\begin{subfigure}[b]{0.4\textwidth}
    	\includegraphics[width=\textwidth]{logo}
    	\caption{More logo.}
  	\end{subfigure}
  	\caption{The same cup of coffee. Two times.}
  	\label{fig:coffee}
\end{figure}
	
Unabhängige Bilder (aber nebeneinander)
	
\begin{figure}[H]
	\centering
	\begin{minipage}{.5\textwidth}
  		\centering
  		\includegraphics[width=\textwidth]{logo}
  		\caption{A figure}
  		\label{fig:test1}
	\end{minipage}%
	\begin{minipage}{.5\textwidth}
  		\centering
  		\includegraphics[width=\textwidth]{logo}
  		\caption{figure}
  		\label{fig:test2}
	\end{minipage}
\end{figure}
	
\subsection{Aufzählungen}\label{subsec:lists}
	
Aufzählungen ohne voranstehende Beschriftung
	
\begin{description}
	\item{Item 1}
	\item{Item 2}
	\item{Item 3}
\end{description}
	
Ungeordnete Aufzählungen
	
\begin{itemize}
	\item{Item 1}
	\item{Item 2}
	\item{Item 3}
\end{itemize}
	
Geordnete Aufzählungen 
	
\begin{enumerate} %[a/1/I/A] optional
	\item{Item 1}
	\begin{enumerate}
		\item Auch verschachtelt möglich
		\item Auch verschachtelt möglich
	\end{enumerate}
	\item{Item 2}
	\item{Item 3}
\end{enumerate}
	
Weitere Umgebungen sind compactenum und compactitem
	
\subsection{Tabellen}\label{subsec:tables}
	
Im folgenden eine Tabelle mit 2 Spalten und mehreren Zeilen sowie einer alternierenden Zeilenfarbe
	
\begin{table}[H]
	\centering
	\caption{Beispieltabelle}
	\label{tab:test}
	\rowcolors{2}{gray!25}{white}
	\begin{tabular}{|c|c|}
    	\rowcolor{gray!50}
    	\hline
    	Table head & Table head\\
    	Some values & Some values\\
    	Some values & Some values\\
    	Some values & Some values\\
    	Some values & Some values\\
    	\hline
  	\end{tabular}
	\sourcecenter{Das ist einfach nur ein Test}
\end{table}
	
\subsection{Abkürzungen}\label{subsec:abbrev}
	
Das ist ein \ac{html}
	
\newpage
	
Das ist immer noch ein \ac{html}. \acp{dr} sind auch ganz toll.
	
\subsection{Source-Code}\label{subsec:code}
	
\begin{lstlisting}[style=htmlcssjs,float=h,caption={JavaScript Hello World Beispiel},label=lst:hello]
var x = "Hello World" //Das ist ein einfacher String
console.log(x.toUppercase());
/* Dieser Code ist sehr cool aber
effektiv bringt er eigentlich nichts. Hähäh */
\end{lstlisting}
	
\autoref{lst:hello} ist ziemlich cool.

\begin{lstlisting}[style=htmlcssjs,float=h,caption={Einfaches HTML Beispiel},label=lst:html]
<!DOCTYPE html>
<html>
	<body class="Hallo">
		<h1>My First Heading</h1>
		<p>My first paragraph.</p>
	</body>
</html>
\end{lstlisting}

\includewebcode{Das ist ein Beispiel}{test.html}

Ich versuche jetzt mal ein autoref auf \autoref{lst:test.html}
	
\end{document}